\RequirePackage{siunitx}  %needed for formatting
\RequirePackage{xparse}   %needed for optional arg type G (printPoly)
\ExplSyntaxOn%
\cs_if_exist:NF \fp_format:nn { \RequirePackage{l3str-format} } % needed for TL2025
% \debug_on:n { check-declarations }

%TODO: eval poly -- use (#1)^expo to avoid sign errors
%TODO: add printWithParens as key to printNum

%%%%%%%%%%%%%%%%%%%%%%%%%%%%%%%%%%%%%%%%%%%%%%%%%%%%%%%%%%%%%%%%%%%%%%%%%%%%%%%%
%% Key definitions
%%%%%%%%%%%%%%%%%%%%%%%%%%%%%%%%%%%%%%%%%%%%%%%%%%%%%%%%%%%%%%%%%%%%%%%%%%%%%%%%

\keys_define:nn { print/coord }
{
  delim    .tl_set:N        = \l__print_coord_delim_tl,
  delim    .initial:n       = {,},
  parens   .code:n          = {
    \tl_set:Nx \l__print_coord_lparens_tl { \tl_item:nn {#1} {1} }
    \tl_set:Nx \l__print_coord_rparens_tl { \tl_item:nn {#1} {2} }
  },
  parens   .initial:n       = {()},
  parens   .value_required:n = true,
  lparens  .tl_set:N        = \l__print_coord_lparens_tl,
  rparens  .tl_set:N        = \l__print_coord_rparens_tl,
}

\keys_define:nn { print/frac }
{
  plus  .bool_set:N  = \l__printfrac_plus_bool,
  plus  .initial:n   = false,
  plus  .default:n   = true,
  d     .bool_set:N  = \l__printfrac_displaystyle_bool,
  d     .initial:n   = false,
  d     .default:n   = true,
}

\keys_define:nn { print/num }
{
  %default: set when called without setting
  format      .tl_set:N         = \l__printcoeff_fmt_tl,
  format      .value_required:n = true,
  fmt         .meta:n           = { format = #1 }, %alias for format
  fmt         .value_required:n = true,
  noplus      .bool_set:N       = \l__printcoeff_noplus_bool,
  noplus      .initial:n        = false,
  noplus      .default:n        = true,
  numDec      .int_set:N        = \l_printcoeff_numDec_int,
  numDec      .initial:n        = 4,
  numDec      .value_required:n = true,
  signs-only  .bool_set:N       = \l__signonly_units_bool,
  signs-only  .initial:n        = true,
  signs-only  .default:n        = false,
  one         .meta:n           = { signs-only = #1 }, %alias for `signs-only'
  one         .initial:n        = true,
  one         .default:n        = false,
  zero        .bool_set:N       = \l__print_zero_bool,
  zero        .initial:n        = false,
  zero        .default:n        = true,
  all         .meta:n           = { signs-only = false, zero = true },
}

\keys_define:nn { print/poly }
{
  deriv   .int_set:N         = \l__printpoly_deriv_int,
  deriv   .initial:n         = 0,
  deriv   .default:n         = 1,
  format  .tl_set:N          = \l__printpoly_fmt_tl,
  format  .value_required:n  = true,
  fmt     .meta:n            = { format = #1 }, %alias for format
  fmt     .value_required:n  = true,
  sort    .tl_set:N          = \l__printpoly_sort_tl,
  sort    .initial:n         = >,
  sort    .value_required:n  = true,
  var     .tl_set:N          = \l__printpoly_var_tl,
  var     .initial:n         = x,
  var     .value_required:n  = true,
}

\keys_define:nn {}
{
  % printNum doesn't use delim or parens, but this
  % inheritance works with the print_coord func
  print/num .inherit:n = print/coord
}

%%%%%%%%%%%%%%%%%%%%%%%%%%%%%%%%%%%%%%%%%%%%%%%%%%%%%%%%%%%%%%%%%%%%%%%%%%%%%%%%
%% Shared variables
%%%%%%%%%%%%%%%%%%%%%%%%%%%%%%%%%%%%%%%%%%%%%%%%%%%%%%%%%%%%%%%%%%%%%%%%%%%%%%%%

\bool_new:N  \l__printFuncs_printSign_bool
\fp_new:N    \l__tmpa_fp
\int_new:N   \l__tmpa_int
\fp_new:N    \l__printpoly_coeff_fp
\fp_new:N    \l__printpoly_expo_fp
\int_new:N   \l__printpoly_expoidx_int
\seq_new:N   \l__printpoly_expo_seq
\tl_new:N    \l__printpoly_coeff_tl
\bool_new:N  \l__printpoly_firstTerm_bool
\bool_new:N  \l__printpoly_expozero_bool
\bool_new:N  \l__printpoly_coeffOne_bool
\tl_new:N    \l__polystr_result_tl

%%%%%%%%%%%%%%%%%%%%%%%%%%%%%%%%%%%%%%%%%%%%%%%%%%%%%%%%%%%%%%%%%%%%%%%%%%%%%%%%
%% Helper: \__poly_build_expo_seq:nn
%%   Builds \l__printpoly_expo_seq from coeffs #1 and raw expo arg #2.
%%   If #2 is blank, auto-generates 0..n-1 and applies sort.
%%%%%%%%%%%%%%%%%%%%%%%%%%%%%%%%%%%%%%%%%%%%%%%%%%%%%%%%%%%%%%%%%%%%%%%%%%%%%%%%

\cs_new_protected:Npn \__poly_build_expo_seq:nn #1 #2
  {
    %TODO: error message if #1 and #2 are different lengths

    % iterate over #1, store #2 as a sequence
    % build sequence if #2 is empty
    \seq_clear:N \l__printpoly_expo_seq
    \tl_if_blank:nTF { #2 }
      {
        %build exponents 0..n-1
        \int_step_inline:nnnn { 0 } { 1 } { \clist_count:n {#1} - 1 }
          { \seq_put_right:Nn \l__printpoly_expo_seq {##1} }
        %logic to reverse order
        \str_case:en { \l__printpoly_sort_tl }
          {
            { < } { }
            { > } { \seq_reverse:N \l__printpoly_expo_seq }
          }
      }
      { \seq_set_from_clist:Nn \l__printpoly_expo_seq { #2 } }
  }

%%%%%%%%%%%%%%%%%%%%%%%%%%%%%%%%%%%%%%%%%%%%%%%%%%%%%%%%%%%%%%%%%%%%%%%%%%%%%%%%
%% Helper: \__poly_apply_deriv:
%%   Applies derivative/integral power rule to \l__printpoly_coeff_fp
%%   and \l__printpoly_expo_fp using \l__printpoly_deriv_int.
%%%%%%%%%%%%%%%%%%%%%%%%%%%%%%%%%%%%%%%%%%%%%%%%%%%%%%%%%%%%%%%%%%%%%%%%%%%%%%%%

\cs_new_protected:Npn \__poly_apply_deriv:
  {
    % derivative (deriv > 0)
    \int_step_inline:nn { \l__printpoly_deriv_int }
      {
        \int_set:Nn \l__tmpa_int { ##1 }
        \fp_set:Nn \l__tmpa_fp { \l__printpoly_expo_fp + 1 - \l__tmpa_int }
        \fp_set:Nn \l__printpoly_coeff_fp { \l__printpoly_coeff_fp * \l__tmpa_fp }
      }
    % integral (deriv < 0)
    \int_step_inline:nn { -\l__printpoly_deriv_int }
      {
        \int_set:Nn \l__tmpa_int { ##1 }
        \fp_set:Nn \l__tmpa_fp { \l__printpoly_expo_fp + 2 - \l__tmpa_int }
        \fp_compare:nNnF { \l__tmpa_fp } = { 0 }
          { \fp_set:Nn \l__printpoly_coeff_fp
              { \l__printpoly_coeff_fp / \l__tmpa_fp } }
      }
    % adjust exponent
    \fp_set:Nn \l__printpoly_expo_fp
      { \l__printpoly_expo_fp - \l__printpoly_deriv_int }
  }

%%%%%%%%%%%%%%%%%%%%%%%%%%%%%%%%%%%%%%%%%%%%%%%%%%%%%%%%%%%%%%%%%%%%%%%%%%%%%%%%
%% Helper: \__poly_load_term:nn
%%   Sets \l__printpoly_coeff_fp, \l__printpoly_expo_fp,
%%   \l__printpoly_expozero_bool, and \l__printpoly_coeffOne_bool
%%   for the current term (coeff #1, index into expo seq already set).
%%   Also applies derivative scaling.
%%%%%%%%%%%%%%%%%%%%%%%%%%%%%%%%%%%%%%%%%%%%%%%%%%%%%%%%%%%%%%%%%%%%%%%%%%%%%%%%

\cs_new_protected:Npn \__poly_load_term:n #1
  {
    % get current coeff and expo
    \fp_set:Nn \l__printpoly_coeff_fp { #1 }
    \fp_set:Nn \l__printpoly_expo_fp
      { \seq_item:Nn \l__printpoly_expo_seq { \l__printpoly_expoidx_int } }

    %deriv/integral (VERY rudimentary) power rule
    \__poly_apply_deriv:
    \bool_set:Nn \l__printpoly_expozero_bool
      { \fp_compare_p:nNn { \l__printpoly_expo_fp } = { 0 } }
    \fp_set:Nn \l__tmpa_fp { abs(\l__printpoly_coeff_fp) }
    \bool_set:Nn \l__printpoly_coeffOne_bool
      { \fp_compare_p:nNn { \l__tmpa_fp } = { 1 } }
  }

%%%%%%%%%%%%%%%%%%%%%%%%%%%%%%%%%%%%%%%%%%%%%%%%%%%%%%%%%%%%%%%%%%%%%%%%%%%%%%%%
%% Helper: \__poly_typeset_varexpo:
%%   Typesets the variable and exponent for the current term.
%%   Used by \print_poly:nnnn.
%%%%%%%%%%%%%%%%%%%%%%%%%%%%%%%%%%%%%%%%%%%%%%%%%%%%%%%%%%%%%%%%%%%%%%%%%%%%%%%%

\cs_new_protected:Npn \__poly_typeset_varexpo:
  {
    \fp_compare:nNnF { \l__printpoly_expo_fp } = { 0 }
      {
        %Print if expo nonzero
        \ensuremath{ \l__printpoly_var_tl }
        %Print expo if not one
        \fp_compare:nNnF { \l__printpoly_expo_fp } = { 1 }
          { \ensuremath{^{ \printNum[noplus,one]{ \l__printpoly_expo_fp } }} }
      }
  }

%%%%%%%%%%%%%%%%%%%%%%%%%%%%%%%%%%%%%%%%%%%%%%%%%%%%%%%%%%%%%%%%%%%%%%%%%%%%%%%%
%% Helper: \__polystr_append_varexpo:
%%   Appends the variable and exponent string for the current term
%%   to \l__polystr_result_tl. Used by \polystr_build:nnn.
%%%%%%%%%%%%%%%%%%%%%%%%%%%%%%%%%%%%%%%%%%%%%%%%%%%%%%%%%%%%%%%%%%%%%%%%%%%%%%%%

\cs_new_protected:Npn \__polystr_append_varexpo:
  {
    \fp_compare:nNnF { \l__printpoly_expo_fp } = { 0 }
      {
        \tl_gput_right:NV \l__polystr_result_tl \l__printpoly_var_tl
        \fp_compare:nNnF { \l__printpoly_expo_fp } = { 1 }
          { \tl_gput_right:Nx \l__polystr_result_tl
              { ^ { \fp_eval:n { \l__printpoly_expo_fp } } } }
      }
  }

%%%%%%%%%%%%%%%%%%%%%%%%%%%%%%%%%%%%%%%%%%%%%%%%%%%%%%%%%%%%%%%%%%%%%%%%%%%%%%%%
% \printNum * [keys] {<coeff>}
% #1 - starred variant suppress leading '+' (equivalent to noplus key)
% #2 - optional args: (default)
%       format=<fmt>: ()      fp_format string (default: empty)
%             noplus: (false) toggle suppressing leading '+' without the star
%             numdec: (4)     number of decimal places to round to
%                one: (true)  toggle printing one (alias for signs-only)
%               zero: (false) toggle printing zero
%%%%%%%%%%%%%%%%%%%%%%%%%%%%%%%%%%%%%%%%%%%%%%%%%%%%%%%%%%%%%%%%%%%%%%%%%%%%%%%%
%\num from siunitx for explicit printing of +
\cs_new_protected:Npn \__printcoeff_num:nn #1 #2
  { \num[#1]{#2} }

\NewDocumentCommand \printNum { s O{} m }
  { \print_coeff:nnn {#1} {#2} {#3} }

\cs_new_protected:Npn \print_coeff:nnn #1 #2 #3
  {
    \group_begin:
      \keys_set:nn { print/num } { #2 }

      % Star is equivalent to noplus
      \bool_if:nT {#1} { \bool_set_true:N \l__printcoeff_noplus_bool }

      % Derive print_sign from noplus
      \bool_set_true:N \l__printFuncs_printSign_bool
      \bool_if:NT \l__printcoeff_noplus_bool
        { \bool_set_false:N \l__printFuncs_printSign_bool }

      \fp_set:Nn \l__tmpa_fp { \fp_eval:n { #3 } } % numeric form for comparisons
      \bool_if:nTF {
        \fp_compare_p:n { \fp_abs:n { \l__tmpa_fp } == 1 }
        && \l__signonly_units_bool
      }{
        \fp_compare:nTF { \l__tmpa_fp == -1 }
          { \ensuremath{-} }     %-1 w/ signonly_units
          {
            \bool_if:NTF \l__printFuncs_printSign_bool
              { \ensuremath{+} } %1 w/  sign
              { }                %1 w/o sign
          }
      }{
        \bool_if:nT {% print zero or non-zero number: generic case
          \l__print_zero_bool || \fp_compare_p:n { \l__tmpa_fp != 0 }
        }{
          \group_begin:
            \ensuremath{
              \exp_args:Ne \__printcoeff_num:nn
                {
                  group-separator={,},
                  group-digits=integer,
                  \bool_if:NT \l__printFuncs_printSign_bool
                    { print-implicit-plus=true }
                }
                {
                  \exp_args:NnV \fp_format:nn
                    { round(\l__tmpa_fp, \l_printcoeff_numDec_int) }
                    \l__printcoeff_fmt_tl
                }
            }
          \group_end:
        }
      }
    \group_end:
  }

%%%%%%%%%%%%%%%%%%%%%%%%%%%%%%%%%%%%%%%%%%%%%%%%%%%%%%%%%%%%%%%%%%%%%%%%%%%%%%%%
% \coord * [keys] {<coeff>}
% #1 - optional args : inherited from printNum
% #2 - a clist of values to print separated by ,
%%%%%%%%%%%%%%%%%%%%%%%%%%%%%%%%%%%%%%%%%%%%%%%%%%%%%%%%%%%%%%%%%%%%%%%%%%%%%%%%

\bool_new:N \l__coord_firstTerm_bool

\NewDocumentCommand \coord { O{} m }
  { \print_coord:nn {#1} {#2} }

\cs_new_protected:Npn \print_coord:nn #1 #2
  {
    \group_begin:
      \keys_set:nn { print/coord } { #1 }
      \bool_set_true:N \l__coord_firstTerm_bool
      \ensuremath{ \l__print_coord_lparens_tl }
      \clist_map_inline:nn { #2 }
        {
          \bool_if:NF \l__coord_firstTerm_bool
            { \ensuremath{ \l__print_coord_delim_tl } }
          \printNum* [all, #1] { ##1 }
          \bool_set_false:N \l__coord_firstTerm_bool
        }
      \ensuremath{ \l__print_coord_rparens_tl }
    \group_end:
  }

%%%%%%%%%%%%%%%%%%%%%%%%%%%%%%%%%%%%%%%%%%%%%%%%%%%%%%%%%%%%%%%%%%%%%%%%%%%%%%%%
% #1 - starred variant prints leading '+'
% #2 - optional args: (default)
%       plus: (false) toggle printing leading '+'
%          d: (false) print with \displaymode
% #3 - numerator
% #4 - denominator
%%%%%%%%%%%%%%%%%%%%%%%%%%%%%%%%%%%%%%%%%%%%%%%%%%%%%%%%%%%%%%%%%%%%%%%%%%%%%%%%

\cs_new_protected:Npn \print_frac:nnnn #1 #2 #3 #4
  {
    \group_begin:
      \keys_set:nn { print/frac } { #2 }
      \bool_if:nT {#1} { \bool_set_true:N \l__printfrac_plus_bool }
      \ensuremath{
        \bool_if:nT { \l__printfrac_displaystyle_bool } { \displaystyle }
        \bool_if:nT { \l__printfrac_plus_bool }
          { \fp_compare:nNnT {#3/#4} > {0} + }
        \xintTeXsignedFrac{ \xintPIrr{ \fpeval{#3} / \fpeval{#4} } }
      }
    \group_end:
  }

\NewDocumentCommand \printFrac { s O{} m m }
  { \print_frac:nnnn {#1} {#2} {#3} {#4} }

%%%%%%%%%%%%%%%%%%%%%%%%%%%%%%%%%%%%%%%%%%%%%%%%%%%%%%%%%%%%%%%%%%%%%%%%%%%%%%%%
% \printPoly * [keys] {<coeff>}{<expo>}
% #1 - starred variant enable leading '+'
% #2 - optional args: (default)
%              deriv: (0) number of derivative using power rule
%       format=<fmt>: ()  fp_format string
%               sort: (>) sort auto-generated expo args
%                var: (x) variable used in expression
% #3 - clist of coefficients
% #4 - (optional) clist of exponents

%% G{} allows completely omitting braces for #4, but requires xparse
%% O{} requires using []
%%%%%%%%%%%%%%%%%%%%%%%%%%%%%%%%%%%%%%%%%%%%%%%%%%%%%%%%%%%%%%%%%%%%%%%%%%%%%%%%

\NewDocumentCommand \printPoly { s O{} m G{} }
  { \print_poly:nnnn {#1} {#2} {#3} {#4} }

\cs_new_protected:Npn \print_poly:nnnn #1 #2 #3 #4
  {
    \group_begin:
      \keys_set:nn { print/poly } { #2 }

      \bool_set_true:N \l__printpoly_firstTerm_bool
      \bool_if:nT {#1} { \bool_set_false:N \l__printpoly_firstTerm_bool }

      % iterate over #3, store #4 as a sequence
      % build sequence if #4 is empty
      \__poly_build_expo_seq:nn { #3 } { #4 }
      \int_zero:N \l__printpoly_expoidx_int
      \clist_map_inline:nn { #3 }
        {
          \int_incr:N \l__printpoly_expoidx_int
          % get current coeff and expo, perform power rule
          \__poly_load_term:n { ##1 }
          \fp_compare:nNnF { \l__printpoly_coeff_fp } = { 0 }
            {
              \tl_set:Nn \l__printpoly_coeff_tl { \l__printpoly_coeff_fp }
              \bool_case:n
                {
                  { \bool_lazy_and_p:nn        % first term, cx
                      { \l__printpoly_firstTerm_bool }
                      { !\l__printpoly_expozero_bool } }
                    { \printNum*{ \l__printpoly_coeff_tl }
                      \bool_set_false:N \l__printpoly_firstTerm_bool }
                  { \bool_lazy_and_p:nn        % first term, c
                      { \l__printpoly_firstTerm_bool }
                      { \l__printpoly_expozero_bool } }
                    { \printNum*[one]{ \l__printpoly_coeff_tl }
                      \bool_set_false:N \l__printpoly_firstTerm_bool }
                  { \bool_lazy_all_p:n {       % not first, \pm 1
                      { !\l__printpoly_firstTerm_bool }
                      {  \l__printpoly_expozero_bool }
                      {  \l__printpoly_coeffOne_bool } } }
                    { \printNum[one]{ \l__printpoly_coeff_tl } }
                  { \bool_lazy_and_p:nn        % not first, \pm c or \pm cx
                      { !\l__printpoly_firstTerm_bool }
                      { !\l__printpoly_expozero_bool ||
                        !\l__printpoly_coeffOne_bool } }
                    { \printNum{ \l__printpoly_coeff_tl } }
                }
              %Print math coeffs and vars
              \__poly_typeset_varexpo:
            }
        }
    \group_end:
  }

%%%%%%%%%%%%%%%%%%%%%%%%%%%%%%%%%%%%%%%%%%%%%%%%%%%%%%%%%%%%%%%%%%%%%%%%%%%%%%%%
% \printPolyStr * [keys] {<coeff>}{<expo>}
% #1 - optional args: (default)
%              deriv: (0) number of derivative using power rule
%       format=<fmt>: ()  fp_format string
%               sort: (>) sort auto-generated expo args
%                var: (x) variable used in expression
% #2 - clist of coefficients
% #3 - (optional) clist of exponents
%%%%%%%%%%%%%%%%%%%%%%%%%%%%%%%%%%%%%%%%%%%%%%%%%%%%%%%%%%%%%%%%%%%%%%%%%%%%%%%%

\NewDocumentCommand \printPolyStr { O{} m G{} }
  {
    \polystr_build:nnn {#1} {#2} {#3}
    \ensuremath{ \tl_use:N \l__polystr_result_tl }
  }

\cs_new_protected:Npn \polystr_build:nnn #1 #2 #3
  {
    \group_begin:
      \keys_set:nn { print/poly } { #1 }

      \__poly_build_expo_seq:nn { #2 } { #3 }
      \tl_gclear:N \l__polystr_result_tl
      \int_zero:N \l__printpoly_expoidx_int

      \clist_map_inline:nn { #2 }
        {
          \int_incr:N \l__printpoly_expoidx_int
          \__poly_load_term:n { ##1 }
          \fp_compare:nNnF { \l__printpoly_coeff_fp } = { 0 }
            {
              % Sign / coefficient
              \tl_if_empty:NTF \l__polystr_result_tl
                {
                  % Leading term - no explicit '+'
                  \bool_if:NTF \l__printpoly_coeffOne_bool
                    {
                      \fp_compare:nNnT { \l__printpoly_coeff_fp } = { -1 }
                        { \tl_gput_right:Nn \l__polystr_result_tl { - } }
                    }
                    { \tl_gput_right:Nx \l__polystr_result_tl
                        { \fp_eval:n { \l__printpoly_coeff_fp } } }
                }
                {
                  % Subsequent terms
                  \fp_compare:nNnTF { \l__printpoly_coeff_fp } > { 0 }
                    { \tl_gput_right:Nn \l__polystr_result_tl { + } }
                    { \tl_gput_right:Nn \l__polystr_result_tl { - } }
                  \bool_if:NTF \l__printpoly_coeffOne_bool
                    {
                      \fp_compare:nNnT { \l__printpoly_expozero_bool } = { 1 }
                        { \tl_gput_right:Nn \l__polystr_result_tl { 1 } }
                    }
                    { \tl_gput_right:Nx \l__polystr_result_tl
                        { \fp_eval:n { abs(\l__printpoly_coeff_fp) } } }
                }
              \__polystr_append_varexpo:
            }
        }
    \group_end:
  }

%%%%%%%%%%%%%%%%%%%%%%%%%%%%%%%%%%%%%%%%%%%%%%%%%%%%%%%%%%%%%%%%%%%%%%%%%%%%%%%%
%      \randpolylongdiv [keys] {<coeff>}{<coeff>}
% \randpolyhornerscheme [keys] {<coeff>}{<coeff>}
% #1 - optional args : inherited from printNum
% #2 - a clist of values to print separated by , (dividend)
% #3 - a clist of values to print separated by , (divisor: (1,\n) for Horner)
%%%%%%%%%%%%%%%%%%%%%%%%%%%%%%%%%%%%%%%%%%%%%%%%%%%%%%%%%%%%%%%%%%%%%%%%%%%%%%%%

\NewDocumentCommand \randpolylongdiv { O{} m m }
  {
    \polystr_build:nnn {#1} {#2} {}
    \tl_set:NV \l_tmpa_tl \l__polystr_result_tl
    \polystr_build:nnn {#1} {#3} {}
    \tl_set:NV \l_tmpb_tl \l__polystr_result_tl
    \exp_args:NVV \polylongdiv \l_tmpa_tl \l_tmpb_tl
  }

\NewDocumentCommand \randpolyhornerscheme { O{} m m }
  {
    \polystr_build:nnn {#1} {#2} {}
    \use:e {
      \exp_not:N \polyhornerscheme
        [ x = { \fp_eval:n
            { -\clist_item:nn {#3}{2} / \clist_item:nn {#3}{1} } } ]
        { \tl_use:N \l__polystr_result_tl }
    }
  }

\ExplSyntaxOff